\section{系统鲁棒性与关键参数敏感性综合分析}

在山地极端洪涝灾害应急救援中，实际作战环境常伴随着电池性能退化、气象信道恶化以及指挥机制变动等不确定性扰动。为全面检验本文构建模型的容错能力与鲁棒性，本章从四个核心维度展开系统性敏感性与鲁棒性深度剖析。

\subsection{电池返航安全余量 $\eta$ 极限承载力与分段阶梯递增机理}

在问题一中，本文验证了 $\eta \in [0.10, 0.35]$ 区间内组批方案的变化。进一步分析表明，返航安全余量实质上刻画了无人机应对复杂山地未知风险的“能量缓冲池”：
\begin{enumerate}
    \item \textbf{弹性缓冲裕度}：在 $\eta \le 0.20$ 的充裕区间内，各机型在绝大部分服务区的实际返航剩余 SOC 均保持在 30\% 以上（全网平均返航 SOC 达 39.81\%）。这意味着即便突发 10\% 范围内的逆风或局部绕飞，现有 18 架次方案仍具备充分的自愈冗余，不致引发强制迫降；
    \item \textbf{高敏感阶梯递增}：当安全门限抬高至 $\eta > 0.30$ 时，不仅轻型机无法起飞，重载 C 型机在 S008 等远端受灾点的安全物理载荷急剧跌破单箱物资重量，引发“批次拆解与运力重组”。总能耗由 59.20 kWh 激增至 78.14 kWh（激增 32.0\%），累计作业时间突破 45400 秒。该现象表明，在灾害应急指挥中，安全系数的设定存在清晰的物理临界上限，盲目提高安全指标将导致后勤保障体系发生质的瘫痪。
\end{enumerate}

\subsection{ALNS 启发式算法随机扰动与求解鲁棒性检验}

针对启发式搜索算法可能存在的初值敏感性与早熟停滞缺陷，本文采用控制变量法，固定问题二的物理约束，对算法随机种子进行了多轮扰动试验（Seed 42, 43, 44）。
试验结果表明：
\begin{itemize}
    \item 在所有独立种子运行下，全网 30 箱首批急救保障物资的硬时限违约箱数始终严格收敛于 \textbf{0}，验证了算法中超大惩罚系数（$10^6$）对硬时限边界的绝对压制能力；
    \item 最佳种子 Seed 42 与备选种子 Seed 44 的总架次数分别为 24 架次与 25 架次（变异系数 $CV < 2.5\%$），完工时刻极差仅为 392.7 秒（相对偏差仅 5.0\%），总能耗极差小于 3.3 kWh。这充分证明了本文设计的“Shaw 空间-时间关联破坏 + Regret-2 深度悔恨修复”算子组合具有卓越的全局寻优导向性与解空间稳定性。
\end{itemize}

\subsection{无线电传播信道附加衰减容限与抗毁性分析}

在强降雨或次生大雾天气下，电磁波空中吸收与地表湿润度变化可能引发额外的射频附加衰落。本文对通信链路的抗毁容限进行了应力测试：
\begin{enumerate}
    \item \textbf{回传链路裕度测试}：本文选定的三大空中中继制高点（西区 731.6m、东南 687.3m、东北 676.9m）与调度中心 O01 的固定网关 G01 之间的实际回传损耗分别为 112.8 dB、115.3 dB 与 116.1 dB。与题面标定的 126.0 dB 回传门限相比，分别保有高达 \textbf{13.2 dB}、\textbf{10.7 dB} 与 \textbf{9.9 dB} 的安全余量。这表明即便恶劣降雨引发额外的大气衰减，中继节点与固定网关之间的回传链路仍保有足够的接收裕度，可保证通信链路正常建立；
    \item \textbf{接入链路动态适应性}：在 325 处通信微时段中，接入中继链路的实际损耗平均值仅为 103.4 dB，优于 116.0 dB 门限达 12.6 dB，表明三维制高点选址方案能够为低空深谷作业提供可靠的通信接入保障。
\end{enumerate}

\subsection{跨区资源共享弹性测算}

在问题四中，完全封闭独立分区导致各组无法复用装备，从而在时序并发叠加时出现设备缺口。为评估跨组错峰调度的可行性，本文进一步计算了空闲窗口期的共享效果：
\begin{equation}
    \Delta t_{idle} = t_{start}^{G02}(T011) - t_{back}^{G01}(T003) = 3120.5\text{ s} - 2110.0\text{ s} = 1010.5\text{ s}
\end{equation}
时序核算显示，G01 组某架 B 型机在完成前序任务后存在 1010.5 秒的待命空闲期。若在组内无飞行任务且电池充饱的窗口期内允许临时借调支援（任务完成后返还原组），则：
\begin{itemize}
    \item $K=2$ 方案下 B 型机与 B 型电池的缺口均可降为 0；
    \item 现有 8 架实体机与 14 组电池即可满足全部配送需求。
\end{itemize}
计算表明，适度引入时空错峰的跨区机动协同机制，能够有效降低绝对隔离带来的设备配置成本。
